% DEFINED:   sec:induction-head, sec:icl, sec:small-model-limits
% RESOLVED:  sec:parallel-axes
% FORWARD:   (none)

\section{Is this an induction head?}\label{sec:induction-head}

The circuit we just read out is worth placing in the wider interpretability
literature, because it is a near-cousin of the most studied small circuit in
transformers: the induction head. The pattern our model implements is not the
canonical induction-head pattern, but it is the same \emph{primitive} deployed
for a different purpose, and the distinction is exactly what makes the comparison
useful.

The induction head, from Olsson and colleagues, is a two-head, two-layer circuit
that implements the rule
\[
\ldots \; A \; B \; \ldots \; A \; \to \; B,
\]
that is, on seeing token $A$ again, predict whatever followed the previous
occurrence of $A$. It decomposes into two pieces. A \textbf{previous-token head}
in the first layer makes every position $t$ copy the embedding of token $t - 1$
into its residual, so that position $t$ ends up encoding "I am here, preceded by
token $X$." An \textbf{induction head} in the second layer then queries, from
position $t$, against earlier positions using that previous-token-shifted
signature, finds the position $i$ where the preceding token matched the current
token, and copies what was at position $i$, which is "the thing that followed the
last $A$."

The transformer-circuits framework of Elhage and colleagues decomposes any such
head into a \textbf{QK circuit}, which determines where the attention spike lands
from the query-key product, and an \textbf{OV circuit}, which determines what
value gets copied once the spike is located. The induction head's QK circuit
matches on previous-token-shifted features; its OV circuit copies the value at
the matched position into the residual stream. Our pointer circuit fits the same
frame exactly, just with different content:

\begin{center}
\begin{tabular}{@{}p{0.3\linewidth}p{0.3\linewidth}p{0.3\linewidth}@{}}
\toprule
 & Induction head & Pointer dereference (ours) \\
\midrule
What layer 1 writes into the residual & The previous token's identity, shifted one position & The address bits, gathered from positions 8 and 9 (and 10) into position 10 \\[1ex]
What layer 2's QK matches on & Previous-token signatures, seeking prior occurrences of the current token & The encoded address, seeking the memory position whose positional encoding matches \\[1ex]
What layer 2's OV copies & The token value at the matched position & The memory bit at the matched position \\[1ex]
Domain & Token identities and sequences of them & Positional encodings and an address \\
\bottomrule
\end{tabular}
\end{center}

Both are two-layer compositions of content-addressable lookups, and the
\emph{mechanism} is identical: attention as soft retrieval, with the query
learning to match the right key. Only the \emph{content} differs, token
identities for the induction head against positional addresses for ours. So the
answer to "is it an induction head?" is: structurally yes, literally no. It is
what one might call the address-lookup head, the same primitive in the same
shape, turned on the purer content-addressable-memory task. Induction heads are
the version of this circuit that the statistics of natural language ask for;
ours is the version a memory-addressing task asks for.

\section{In-context learning is content-addressable lookup, composed}\label{sec:icl}

The reason the induction head matters beyond its own neatness is its connection
to \textbf{in-context learning}: the ability of a trained language model to use
information in its prompt, with no weight update, to inform predictions on a task
specified in that prompt. Show a model the prompt

\begin{verbatim}
zarp -> bumi
flom -> bumi
trel -> tev
clok -> tev
zarp -> ?
\end{verbatim}

and it will often emit \texttt{bumi}. The model has, in context, learned a
mapping from a single example with no gradient step. That capability is much of
what separates large language models from the older sequence models, and it is
the substance of the in-context-learning literature.

Olsson and colleagues argue that induction heads are a \emph{mechanistic}
substrate for this ability: a circuit that finds prior occurrences of the current
token and copies the following value is precisely what one-shot
retrieval-from-context needs. Empirically, the emergence of induction heads
during pretraining coincides with the emergence of measurable in-context-learning
ability.

Our pointer-lookup circuit is in the same family. Read the pointer task as a
one-shot in-context-learning problem. The \textbf{context} is an eight-bit memory
and a three-bit address. The \textbf{task specification} is implicit: at this
position, retrieve memory bit $a$. The \textbf{completion} is the value at the
addressed cell. The model was never trained on the specific eleven-bit sequence
it sees at inference time, and most of the twenty thousand training examples
never reappear; it solves each new example using the in-context memory and
address, exactly as an in-context-learning system would. It just happens to do so
by positional-address matching rather than by token-pattern matching, which is
the same primitive applied to a slightly cleaner data-generating process.

This is the broader picture worth carrying away. In-context learning is what a
transformer does when its content-addressable lookup primitive is composed over
enough tokens, heads, and layers to encode arbitrary soft retrieval queries.
Induction heads are one well-studied two-layer instance of that primitive, and
our pointer circuit is another. A trained language model has thousands of
attention heads across dozens of layers, each implementing some specialized
lookup, and in-context learning is the emergent property of running them all
together on a context-dependent task.

Looking back at the inductive-bias frame, this sharpens what a transformer's
architectural commitment actually is. It is not just "attention." It is
specifically content-addressable retrieval as a composable primitive. Depth is
how many composed retrievals the model can chain. Width, the multi-head axis of
\Cref{sec:parallel-axes}, is how many parallel retrievals it runs per layer.
In-context learning emerges where depth and width are large enough that the
composed retrievals can express arbitrary soft programs over the context.

\section{What this small model leaves out}\label{sec:small-model-limits}

A word of honesty about scale before we test that claim. The model interpreted
here is two layers, one head, eleven tokens of context, a two-token vocabulary,
and the circuit is correspondingly minimal. Three things a full interpretability
analysis on a larger model would also have to cover are worth naming, both to be
candid and because the last section runs into the third of them head-on.

\begin{itemize}
\item \textbf{Value circuits and embedding geometry.} We have not opened up the
embedding vectors for tokens 0 and 1 or the positional encodings. With a
two-token vocabulary the embedding has only two rows to inspect, but how
\emph{value} information flows through the OV path of each layer is genuinely
interesting at scale, and we leave it closed here.

\item \textbf{Multi-head specialization.} Real transformers have many heads per
layer, each potentially a different lookup. A character-level language model
trained on prose, for instance, develops heads that attend to the previous space,
heads that attend to the start of the current word, heads that attend to repeated
characters, each a small specialized lookup that together enable rich behavior.
Our one-head model has none of this structure to find. (These per-head
specializations are exactly what mechanistic-interpretability studies of trained
character-level language models visualize; our pointer task is deliberately too
small to grow them.)

\item \textbf{Activation patching.} The ablations here zero out an attention
layer wholesale. A finer tool is to \emph{patch} an activation from one input
into another: copy a residual at one layer and position from one run into another
and watch whether the output flips. This isolates the causal role of each
component on each example, and it is the basis for most published mechanistic
results on language models. The final section leans on exactly this tool, and it
turns out to be the difference between seeing the mechanism and not.
\end{itemize}

The pedagogical point is that the technique applied in this chapter, state a
circuit hypothesis from the architectural argument and verify it with attention
extraction plus ablation, is the same technique used at scale, only with more
components and sharper intervention tools. The next section puts that to the
test. It takes the \emph{same task} at $M = 32$, where \Cref{ch:transformer}
found a third layer is required, and asks whether the clean circuit survives the
scale-up.
